\section{A restricted obstruction for finite moment heads}
\label{sec:obstruction}

The integer hierarchy of \cite[Section 2, Integer moment evolution]{BakerExterior2026}
is
\begin{align}
 S_n'={}&2n\sum_{k=1}^n S_kS_{n+1-k}
    +2n\sum_{k=2}^n\bigl(T_kT_{n+2-k}+DS_kS_{n+2-k}\bigr)
    \notag\\
    &-4n(n+2)S_{n+1}-4n(n+3)DS_{n+2},\qquad n\ge1,
    \label{eq:Shierarchy}\\
 T_n'={}&2(2n-1)\sum_{k=2}^nT_kS_{n+1-k}
    +4nD\sum_{k=2}^nT_kS_{n+2-k}
    \notag\\
    &-4(n^2+n-1)T_{n+1}-4n(n+3)DT_{n+2},\qquad n\ge2.
    \label{eq:Thierarchy}
\end{align}
Empty sums vanish. The nonlinear terms in the $n$th row use indices
at most $n$; the two higher indices occur linearly. This feature, with
the pair law, is the only input needed for the obstruction below.

For clarity about the inherited physical realization, at a finite
simple polynomial slice one may temporarily use
$P_1=\sum_j w_j/(w_j^2-D)$. The exterior velocities relative to the
pair mean are
\begin{equation}
 w_j'=4w_jm_j+2P_1+2\sum_{k\ne j}(w_j-w_k)^{-1},
 \qquad m_j=(w_j^2-D)^{-1}.
 \label{eq:relative}
\end{equation}
Differentiate the weights in \eqref{eq:moments}, pair the last sums
using divided differences, and subtract their diagonals. If
$A=(u^2-D)^{-1}$, $C=(v^2-D)^{-1}$ and
$L_h=\sum_{r=0}^{h-1}A^{h-1-r}C^r$ with $L_0=0$, the identities are
\begin{equation}
 \frac{A^h-C^h}{u-v}=-(u+v)AC L_h,\qquad
 \frac{uA^h-vC^h}{u-v}=-AC L_{h-1}-(D+uv)AC L_h.
 \label{eq:divdiff}
\end{equation}
The endpoints of the convolutions cancel every occurrence of $P_1$,
leaving \eqref{eq:Shierarchy}--\eqref{eq:Thierarchy}. These rational
divided differences also remain regular when two exterior nodes
coincide away from the target collar. The parent's proof supplies the
complete canonical passage using cleared-contour time jets of the
opposite-paired polynomial cutoffs and the inverse-square tail; it
does not assume absolute convergence of $P_1$ or arbitrary-sequence
well-posedness. Those hypotheses remain in force here.

\begin{theorem}[No nonconstant universal finite-head integral]
\label{thm:nohead}
Fix $P\ge1$. Let
\[
 I=I(t,D,S_1,\ldots,S_P,T_2,\ldots,T_P)
\]
be $C^1$ on an open real coordinate domain with $D\ne0$. Suppose
the chain-rule derivative of $I$ vanishes under
\eqref{eq:pair}, \eqref{eq:Shierarchy} and \eqref{eq:Thierarchy}
for freely variable tails through order $P+2$. Then $I$ is locally
constant. Here ``freely variable'' means that the relevant head and
tail coordinates range through an open set, with no extra relation
fixing the higher moments from the retained ones.
\end{theorem}
\begin{proof}
In $\dd I/\dd t$, the coefficients of $S_{P+2}$ and $T_{P+2}$
are respectively
\begin{equation}
 -4P(P+3)D I_{S_P},\qquad -4P(P+3)D I_{T_P}.
 \label{eq:highest-tail}
\end{equation}
The odd term is absent when $P=1$. No lower row contains either
highest coordinate, and the nonlinear terms contain neither.
Independent variation and $D\ne0$ imply that these partial
derivatives vanish. Locally remove those variables from $I$ and
repeat the argument at $P-1$. Eventually $I=I(t,D)$ and its
remaining equation is
\[
 I_t+(2-4DM)I_D=0.
\]
Variation of $M$ now forces $I_D=0$ and $I_t=0$. All first
derivatives of $I$ vanish on the domain, so it is locally constant.
\end{proof}

A $C^1$ first integral on an open neighborhood crossing $D=0$ cannot
acquire a nonconstant value only on that hypersurface: its derivatives
vanish on the dense complement and hence there too. This does not
exclude functions defined solely on the hypersurface, which is not an
invariant flow domain since $D'=2$ there.

\begin{proposition}[A finite polynomial realization of the free-tail premise]
\label{prop:rank}
At fixed real $D\ne0$, take $2P+3$ freely located real exterior
roots in a strict collar complement. There exist configurations where
the map from these roots to
\[
 (S_1,\ldots,S_{P+2},T_2,\ldots,T_{P+2})
\]
has full rank $2P+3$. Full rank persists on an open set and is generic
in any connected analytic configuration region on which the resulting
determinant is not identically zero.
\end{proposition}
\begin{proof}
The moment kernels are $f_r(w)=(w^2-D)^{-r}$ for $1\le r\le P+2$
and $g_r(w)=w(w^2-D)^{-r}$ for $2\le r\le P+2$. At infinity,
\[
 f_r'(w)=-2r w^{-2r-1}+O(w^{-2r-3}),\qquad
 g_r'(w)=(1-2r)w^{-2r}+O(w^{-2r-2}).
\]
Their leading powers are all distinct. In any putative nontrivial
linear relation, the least decaying nonzero term could not cancel.
They are consequently linearly independent rational functions.
Their evaluation vectors on an open real half-line outside the collar
span $\R^{2P+3}$: otherwise one nonzero fixed linear functional
would annihilate all evaluations, giving a rational identity there
and everywhere. Select distinct points whose evaluation vectors form
a basis. The moment Jacobian at those roots is precisely this
nonsingular evaluation matrix. Continuity and analyticity give the
stated open-set and generic conclusions.
\end{proof}

Each such initial divisor, together with the target pair, is a genuine
polynomial initial slice and has its finite heat evolution. Thus the
obstruction is not solely about fictitious independent sequences.
It remains a statement about a freely variable polynomial exterior
class. It does not establish free tails for the unchanged theta
source, exclude an invariant along one prescribed canonical
trajectory, treat an infinite functional, or cover a fixed small
divisor whose moments satisfy algebraic relations. In particular it
does not contradict Section~\ref{sec:finite}.
